\begin{table}[h]
\centering

\caption{两点分布、二项分布、超几何分布、几何分布的对比：一个盒子中有\(3\) 个红球、\(7\) 个白球}
\begin{tabular}{>{\centering\arraybackslash}p{2.8cm}
                >{\centering\arraybackslash}p{2cm}
                >{\centering\arraybackslash}p{4cm}
                >{\centering\arraybackslash}p{6cm}}
\toprule
\textbf{分布} & \textbf{是否放回} & \textbf{试验/摸球次数} & \textbf{随机变量 \(X\) 的含义} \\
\midrule
两点分布      & 不涉及 & 固定为 \(1\) 次      & 成功（红球）与否（\(X=1\) 或 \(0\)） \\
\hline
二项分布      & 有放回           & 固定（例如 \(n=5\) 次） & 成功（红球）的总次数 \\
\hline
超几何分布    & 不放回           & 固定（例如 \(n=5\) 次） & 成功（红球）的总次数 \\
\hline
几何分布      & 有放回           & 不固定（直到首次成功） & 首次成功时的摸球次数 \\
\bottomrule
\end{tabular}
\end{table}

\begin{table}[h]
\centering
\caption{两点分布、二项分布、超几何分布、几何分布的期望与方差}
\begin{tabular}{>{\centering\arraybackslash}p{3cm}
                >{\centering\arraybackslash}p{3.5cm}
                >{\centering\arraybackslash}p{4cm}}
\toprule
\textbf{分布} & \textbf{均值 \(E(X)\)} & \textbf{方差 \(D(X)\)} \\
\midrule
两点分布 \(X \sim \text{Bernoulli}(p)\) 
& \(p\) 
& \(p(1-p)\) \\
\hline
二项分布 \(X \sim B(n, p)\) 
& \(np\) 
& \(np(1-p)\) \\
\hline
超几何分布 \(X \sim H(n, M, N)\) 
& \(n \cdot \dfrac{M}{N}\) 
& \(n \cdot \dfrac{M}{N} \cdot \left(1-\dfrac{M}{N}\right) \cdot \dfrac{N-n}{N-1}\) \\
\hline
几何分布 \(X \sim \text{Geometric}(p)\) 
& \(\dfrac{1}{p}\) 
& \(\dfrac{1-p}{p^{2}}\) \\
\bottomrule
\end{tabular}
\end{table}

\subsection*{参数说明}
\begin{itemize}
  \item 两点分布（伯努利分布）：每次试验成功概率为 \(p\)，\(X=1\) 表示成功，\(X=0\) 表示失败。
  \item 二项分布：\(n\) 次独立重复伯努利试验，每次成功概率 \(p\)。
  \item 超几何分布：总体大小 \(N\)，其中成功个体数 \(M\)，不放回抽取 \(n\) 个，\(X\) 为成功个数。
  \item 几何分布：每次试验成功概率 \(p\)，\(X\) 为首次成功时的总试验次数。
\end{itemize}

\clearpage

设 $X \sim \text{Geometric}(p)$，表示首次成功所需的总试验次数，每次试验成功概率为 $p$，失败概率为 $q = 1-p$。分布列为：
$
P(X=k) = q^{k-1} p, \quad k=1,2,3,\dots
$

\subsection*{第一步：利用无记忆性建立递推}

考虑第一次试验的结果：
\begin{itemize}
    \item 以概率 $p$：第一次成功，此时 $X=1$。
    \item 以概率 $q$：第一次失败，此时局面“重置”，但已消耗 1 次试验，后续还需 $X'$ 次（$X'$ 与 $X$ 同分布且独立），故 $X = 1 + X'$。
\end{itemize}

\subsection*{第二步：推导 $E[X]$ }
$
E[X] = p\cdot1 + q\cdot(1+E[X]) \quad\Rightarrow\quad E[X]=\dfrac{1}{p}
$

\subsection*{第三步：推导 $E[X^2]$}
$
E[X^2] = p\cdot1^2 + q\cdot E[(1+X)^2] = p + q\bigl(1+2E[X]+E[X^2]\bigr)
$

代入 $E[X]=\dfrac{1}{p}$：
$
E[X^2] = 1 + \dfrac{2q}{p} + qE[X^2] \quad\Rightarrow\quad pE[X^2] = 1 + \dfrac{2q}{p}
$

$
E[X^2] = \dfrac{1}{p} + \dfrac{2q}{p^2} = \dfrac{p+2q}{p^2} = \dfrac{2-p}{p^2}
$

\subsection*{第四步：计算方差}
$
D(X) = E[X^2] - (E[X])^2 = \dfrac{2-p}{p^2} - \dfrac{1}{p^2} = \dfrac{1-p}{p^2}
$
